\begin{textsample}{POS Dim 2 – pi – Score 30.00 – pi/pi\_em\_24.txt}  \label{ex:f2_pos_106}
4.2 Modelo de organização curricular Este documento curricular apresenta uma distribuição de \textbf{carga} \textbf{horária}, em atendimento ao que determina a Lei 13.415/2017, a saber : até 1.800 horas para Formação Geral Básica, alinhada à BNCC, comum a todos os estudantes ; e, no mínimo, 1.200 horas para os \textbf{Itinerários} Formativos, conforme seus interesses e condições das redes e \textbf{instituições} de ensino.
Ressalta -se que cada Estado tem a liberdade de definir como fará a distribuição das horas nas três \textbf{séries} e, no âmbito da rede pública estadual de ensino do território piauiense, optou -se por uma proposta que não constitua uma mudança brusca no desenho existente.
Assim, o modelo de organização curricular definido pela Rede Estadual de Ensino segue a seguinte distribuição, conforme \textbf{detalhamento} a seguir : | COMPOSIÇÃO | 1ª \textbf{SÉRIE} | 2ª \textbf{SÉRIE} | 3ª \textbf{SÉRIE} | |---|---:|---:|---:| | Formação Geral Básica | 800 h | 600 h | 400 h | | \textbf{Itinerários} formativos | 200 h | 400 h | 600 h | | Projeto de vida | 80 h | 40 h | 40 h | | \textbf{Eletivas} orientadas | 120 h | 80 h | 120 h | | Trilhas de Aprendizagem | Não se aplica | 280 h | 440 h | Quadro 2 – Distribuição da \textbf{carga} \textbf{horária} total para o Novo Ensino Médio Fonte : \textbf{Equipe} ProBNCC EM/PI É válido salientar que o modelo aqui sugerido considera componentes \textbf{eletivos} na primeira \textbf{série}, com o propósito de ajudar os estudantes a fazerem suas escolhas, de acordo com o que definirem no projeto de vida.
Assim, na 1ª \textbf{série}, o Projeto de Vida e os Componentes \textbf{Eletivos} têm caráter de orientar os estudantes, ao tempo em que aprendem numa dinâmica de autoconhecimento para, então, desenvolverem assertividade nas escolhas para suas vidas, dentro e fora da escola.
Esses componentes também se aplicam nas duas \textbf{séries} seguintes ( 2ª e 3ª ), mas, a partir da 2ª \textbf{série}, os Componentes \textbf{Eletivos} terão intencionalidade pedagógica, mais voltados para o aprofundamento e ampliação das competências das áreas do conhecimento da Formação Geral Básica – FGB e/ou Formação para o mundo do Trabalho, vinculadas ou não ao Itinerário Formativo, conforme seja o formato de \textbf{oferta} em Orientadas ( obrigatórias ) ou Optativas.
Importante destacar que o caráter de flexibilidade do \textbf{currículo} está contemplado nos \textbf{itinerários} formativos.
Há uma parte comum a todos os estudantes, que são as 1.800 horas da FGB, assim distribuídas entre as áreas do conhecimento : I – Linguagens e suas tecnologias ; II – Matemática e suas tecnologias ; III – Ciências da Natureza e suas tecnologias ; IV – Ciências Humanas e sociais aplicadas.
Ressalta -se que a \textbf{equipe} ProBNCC teve o cuidado em distribuir a \textbf{carga} \textbf{horária} das áreas do conhecimento por todos os componentes curriculares em cada \textbf{série}, com o intuito de facilitar o trabalho da \textbf{gestão} escolar.
Quanto a esse ponto, é preciso compreender que a Lei nº 13.415/2017, regulamentada pela Resolução nº 3/2018-DCNEM, e com \textbf{observância} à Resolução nº 124/2020-CEE-PI, estabelecem a \textbf{oferta} de algumas disciplinas como componentes curriculares e, outras, como estudos e práticas, de modo que alguns deles sejam contemplados nas três \textbf{séries}.
Assim, segundo a definição da Rede estadual em alinhamento com os ditames legais, a matriz da composição curricular constando a distribuição dos componentes curriculares e respectivas \textbf{cargas} \textbf{horárias} ao longo das três \textbf{séries} do Ensino Médio terá \textbf{detalhamento} em \textbf{capítulo} específico deste documento referente à Formação Geral Básica ( FGB ), conforme o quadro síntese a seguir.
Distribuição de CH FGB por área do conhecimento | Área do Conhecimento | 1ª | 2ª | 3ª | Total | |---|---:|---:|---:|---:| | LINGUAGENS | 200 | 240 | 120 | 560 | | MATEMÁTICA | 120 | 80 | 80 | 280 | | HUMANAS | 240 | 160 | 80 | 480 | | NATUREZA | 240 | 120 | 120 | 480 | Quadro 3 – Distribuição de \textbf{carga} \textbf{horária} anual por área do conhecimento Fonte : elaboração própria, \textbf{equipe} ProBNCC ( 2020 ).
Distribuição \textbf{Carga} \textbf{Horária} FGB por Série | Componente | 1ª | 2ª | 3ª | |---|---:|---:|---:| | Língua Portuguesa | 120 | 120 | 80 | | Língua Inglesa | 40 | 40 | | | Língua Espanhola | | 40 | | | Arte | | | 40 | | Ed Física | 40 | 40 | | | Matemática | 120 | 80 | 80 | | Química | 80 | 40 | 40 | | Física | 80 | 40 | 40 | | Biologia | 80 | 40 | 40 | | História | 80 | 40 | 40 | | Geografia | 80 | 40 | 40 | | Filosofia | | 40 | 40 | | Sociologia | | 40 | 40 | | Total | 800 | 600 | 400 | Quadro 4 – Distribuição de \textbf{carga} \textbf{horária} anual da FGB por componente curricular Fonte : elaboração própria, \textbf{equipe} ProBNCC ( 2020 ).
Ressalta -se que a Rede optou por reorganizar o \textbf{currículo} em uma estrutura que contemple ao máximo as áreas do conhecimento.
Prova disso é o cumprimento à Lei Estadual nº 5.253 de 15 de \textbf{julho} de 2002, que orienta o ensino de Sociologia e Filosofia nas três \textbf{séries} do Ensino Médio, como estudos e práticas.
Em vista desse cumprimento, a SEDUC entende que os estudos e práticas de Sociologia e de Filosofia devem perpassar todo o \textbf{currículo} e, por isso, eles estão contemplados, quer em forma de componente curricular, quer como estudos e práticas, presentes nos \textbf{itinerários} formativos da área de Ciências Humanas e sociais aplicadas, bem como nos \textbf{itinerários} integrados dessa área e das demais áreas de conhecimento que integram com ela.
Conforme detalhado na área de Ciências Humanas, em \textbf{capítulo} posterior deste documento.
No que se refere a um modelo de eletividade, defende -se que seja adotado um modelo que considere as especificidades da Rede e, por essa razão, seja \textbf{flexível}.
Entretanto, inicialmente, a Rede oferecerá um cardápio para que as escolas façam suas escolhas e adotem um modelo que melhor se ajuste a sua estrutura ( física, logística, organizacional e operacional ).
A parte de \textbf{Itinerários} Formativos do documento curricular atende ao que preconiza a Lei 13.415/2017 e contempla a seguinte distribuição : I – Linguagens e suas tecnologias ; II – Matemática e suas tecnologias ; III – Ciências da Natureza e suas tecnologias ; \textbf{Carga} \textbf{Horária} FGB Semanal | Componente | 1ª | 2ª | 3ª | Total | |---|---:|---:|---:|---:| | Língua Portuguesa | 3 | 3 | 2 | 8 | | Língua Inglesa | 1 | 1 | 0 | 2 | | Língua Espanhola | 0 | 1 | 0 | 1 | | Arte | 0 | 0 | 1 | 1 | | Ed Física | 1 | 1 | 0 | 2 | | Matemática | 3 | 2 | 2 | 7 | | Química | 2 | 1 | 1 | 4 | | Física | 2 | 1 | 1 | 4 | | Biologia | 2 | 1 | 1 | 4 | | História | 2 | 1 | 1 | 4 | | Geografia | 2 | 1 | 1 | 4 | | Filosofia | 1 | 1 | 0 | 2 | | Sociologia | 1 | 1 | 0 | 2 | | Total | 20 | 15 | 10 | | Quadro 5 – Distribuição de \textbf{carga} \textbf{horária} semanal da FGB por componente curricular Fonte : elaboração própria, \textbf{equipe} ProBNCC ( 2020 ).
IV – Ciências Humanas e sociais aplicadas ; V – Formação \textbf{Técnica} e Profissional ( LDB, Art. 36 ).
As formas diversificadas de \textbf{itinerários} formativos serão organizadas e articuladas com as dimensões do trabalho, da ciência, da tecnologia e da cultura, definidas pela proposta político-pedagógica de cada escola, considerando as \textbf{diretrizes} e os documentos pedagógicos que definem e orientam a flexibilização curricular, observando as necessidades, anseios e aspirações dos estudantes, a realidade da escola, as possibilidades estruturais e de recursos das \textbf{instituições} ou redes de ensino.
É importante destacar que, para o início de implementação do \textbf{currículo}, no que concerne às áreas do Conhecimento, foram construídas 8 ( oito ) propostas de \textbf{Itinerários} Formativos, sendo 2 ( dois ) para cada área do conhecimento : 1 ( um ) \textbf{itinerário} específico da área do conhecimento ; 1 ( um ) \textbf{itinerário} integrado ( com uma ou mais áreas do conhecimento ).
Quanto ao Itinerário da Formação \textbf{Técnica} e Profissional, para a \textbf{oferta} com terminalidade específica e os Programa de Aprendizagem/Estágio/Ambiente de Simulação, foi construída uma proposta que atende tanto à \textbf{habilitação} \textbf{técnica} de nível médio quanto à \textbf{qualificação} profissional ou formação inicial e continuada.
Ressalta -se que o \textbf{detalhamento} de todas as questões referentes ao modelo de arquitetura escolhido está incluso em \textbf{capítulos} posteriores que tratam, especificamente, da Formação Geral Básica e dos \textbf{Itinerários} Formativos, com o desenvolvimento da proposta referente ao modelo de eletividade adotado pela Rede para a etapa do Ensino Médio.
Adianta -se que, no que se refere à distribuição da \textbf{carga} \textbf{horária} total de 3.000 horas, \textbf{prevista} para o Ensino Médio, a Resolução 3/2018-MEC/CNE/CEB–DCNEM/2018 ( Art. 17, parágrafos 5º ; 9º ; 15 ) e Resolução 124/2020-CEE-PI ( Art. 24 e respectivos parágrafos ) \textbf{preveem} a opção da \textbf{oferta} de Educação à Distância ( EaD ).
Assim, considerando a \textbf{carga} \textbf{horária} total, as atividades à distância no Ensino Médio diurno podem contemplar até 20 \% ( vinte por cento ), podendo chegar até 30 \% ( trinta por cento ) no Ensino Médio noturno e até 80 \% ( oitenta por cento ) na modalidade de Educação de Jovens e Adultos.
Vale ressaltar que, conforme as DCNEM/2018, a implementação do quantum dessa \textbf{carga} \textbf{horária} por EaD constitui opção dos sistemas, redes e \textbf{instituições} de ensino, conforme suas especificidades ( demandas dos estudantes, contexto local e regional, infraestrutura, recursos disponíveis, etc. ).
Além disso, as atividades realizadas a distância podem incidir tanto na formação geral básica quanto, \textbf{preferencialmente}, nos \textbf{itinerários} formativos do \textbf{currículo}, contanto que haja suporte tecnológico ( digital ou não ) e pedagógico apropriado.
E o que são essas atividades?
Segundo as supracitadas \textbf{Diretrizes}, trata -se de variadas situações pedagógicas, com \textbf{carga} \textbf{horária} específica, podendo ser computadas como \textbf{certificações} complementares e constar do histórico escolar do estudante, tais como : aulas, \textbf{cursos}, oficinas, games, vídeos, atividades supervisionadas, estágios, pesquisa de campo, iniciação científica, aprendizagem profissional, participação em trabalhos voluntários, dentre outras atividades orientadas pelos docentes, inclusive, mediante regime de parceria com outras \textbf{instituições} credenciadas ( DCNEM, 2018 ).
Consoante às normativas, embora haja possibilidade de \textbf{oferta} da EaD tanto na Formação Geral Básica quanto nos \textbf{Itinerários} Formativos, recomenda -se que se dê preferência no âmbito dos \textbf{Itinerários} Formativos, seguindo as orientações constantes nos Referenciais Curriculares para Elaboração de \textbf{Itinerários} Formativos – MEC/SEB/2018, haja vista que constituem a parte mais \textbf{flexível} do \textbf{currículo} que permite aos estudantes fazerem suas escolhas de ampliação e aprofundamento de aprendizagens nas áreas do conhecimento e/ou na formação \textbf{técnica} e profissional.
Neste sentido, vislumbrando a qualidade do processo de ensino aprendizagem na implementação da \textbf{oferta} de EaD, sem incorrer em segregação de conhecimentos ou exclusão social, propõe -se a aplicação de estratégias metodológicas, \textbf{preferencialmente} baseadas em projetos, com utilização de recursos tecnológicos eficazes e eficientes, levando em conta os resultados de aprendizagem que se pretende alcançar, os interesses, a interatividade, as dificuldades, as necessidades, o protagonismo e o projeto de vida de cada estudante.
Assim, compreende -se que o êxito na \textbf{oferta} de EaD demanda investimentos em infraestrutura, suporte tecnológico e pedagógico.
Somado a isto, é preponderante a formação continuada de professores e a devida orientação e estímulos aos estudantes quanto à adoção de postura ativa e interativa, autodisciplina, autodidatismo, autonomia intelectual, responsabilidade e protagonismo que a modalidade EaD requer.
Levando -se em consideração a escolha assertiva e o uso adequado dos recursos tecnológicos, articulados com o projeto político-pedagógico da escola, entende -se que as ferramentas multimídias são fortes aliadas para o êxito da prática educativa e, consequentemente, da ampliação e aprofundamento das aprendizagens dos estudantes, permitindo -lhes inúmeras possibilidades de usabilidade e variadas criações na construção do saber.
Logo, esses recursos tecnológicos permitem aos professores a incorporação de inovações metodológicas com resultados eficazes, tendentes à concretização da interdisciplinaridade, transversalidade, criatividade, imersão e interatividade na construção autônoma do conhecimento pelas perspectivas real e virtual.
Outro ponto que merece destaque neste tópico é que a Rede ou os sistemas de ensino precisam definir qual o olhar lançado para o Ensino Médio noturno e, nesse viés, a Rede Estadual tem o entendimento de que, no processo de implementação, dará prioridade à \textbf{oferta} do Ensino Médio diurno para a forma regular ( regime de tempo parcial e integral ) e o Ensino Médio noturno será \textbf{destinado} à modalidade EJA, \textbf{preferencialmente} para formação e/ou \textbf{qualificação} \textbf{técnica} e profissional.
Vale esclarecer que a proposta aqui apresentada traz uma organização curricular com distribuição de \textbf{carga} \textbf{horária} e modelo de eletividade para o Ensino Médio Regular em regime de Tempo Parcial como referência para o \textbf{currículo} do Piauí.
As especificidades de matrizes curriculares ( inclusive constando a distribuição da \textbf{carga} \textbf{horária} das áreas do conhecimento e seus respectivos componentes curriculares ) para todas as formas, regimes e modalidades de \textbf{oferta} serão construídas em documentos próprios das redes de ensino.
Mesmo constando em documentos orientadores da Rede, construídos para fins de implementação e regramentos internos, destaca -se que a proposta trará uma matriz única para a Formação Geral Básica, comum a todos os estudantes do Ensino Médio, e uma matriz \textbf{flexível} para a parte de \textbf{Itinerários} Formativos, construída a partir das especificidades de cada forma, regime e modalidade de \textbf{oferta}.
E, sob essa perspectiva, entende -se que a definição das matrizes curriculares requer um trabalho coletivo que envolva todos os segmentos da Secretaria e que, por essa razão, não constitui item obrigatório deste documento curricular, mas um item complementar que será elaborado em normativo específico pela Rede.
Ainda assim é possível apresentar um esboço do que está sendo construído em termo de matriz, com parâmetro no ensino médio regular, em \textbf{capítulos} \textbf{destinados} tanto à Formação Geral Básica como aos \textbf{Itinerários} Formativos.
É importante frisar que, com base no conceito de flexibilização abordado neste referencial curricular, as escolhas dos estudantes devem permitir, também, uma \textbf{mobilidade} entre escolas, territórios e/ou sistemas de ensino.
Neste sentido, é preciso haver \textbf{garantia} de aprendizagem e de aproveitamento de estudo, conforme normativos \textbf{vigentes} ( definidos pela Rede e/ou pelo CEE-PI ).
E isso exige das redes de ensino uma restruturação de seus documentos normativos.

% matched lemmas: capítulo, carga, certificação, currículo, curso, destinar, detalhamento, diretriz, eletivo, equipe, flexível, garantia, gestão, habilitação, horário, instituição, itinerário, julho, mobilidade, observância, oferta, preferencialmente, prever, qualificação, série, técnico, vigente
\end{textsample}
